% 中文摘要
\begin{abstractzh}
大语言模型驱动的智能体已从单轮问答走向多步工具调用与环境交互，并在部署后通过周期性强化学习持续更新策略。这类训练面临一个现实瓶颈：\textbf{强化学习需要海量交互轨迹与可信奖励，而人工标注既昂贵又无法规模化}。若奖励只依据智能体的自我陈述，模型极易学会"谎报完成"这一奖励作弊行为；若依赖人工评分，则无法支撑持续、大规模的策略更新。

本文提出并实现一套\textbf{面向自进化的智能体强化学习系统}：部署中的智能体在沙箱环境中与工具交互、自主产生经验，由多智能体闭环提供可信的自动奖励，强化学习据此更新策略——模型在无人工标注的条件下自我进化。核心贡献有三：

其一，\textbf{多智能体采样架构}。执行者（actor）在代码沙箱中以 ReAct 方式多步调用工具求解任务；观察者（observer）以沙箱执行前后的\textbf{状态差分}为客观依据生成状态报告；提问者（questioner）扮演具人设的模拟用户，驱动多轮追问，使单个种子任务扩展为一段多轮对话。三者协同，把"一次采样"变成"一段自主交互经验"。

其二，\textbf{差分驱动的可信奖励}。奖励不取自智能体自述，而由一个外部冻结的大模型裁判，依据观察者提供的\textbf{环境真实状态差分}统一打分，采用完成度、正确性、轨迹质量、安全四维度乘性聚合、安全为二元门控。以真实状态差分而非自我陈述为锚点，从机制上抑制奖励作弊。

其三，\textbf{系统实现与端到端验证}。将上述采样与奖励机制整合为一套完整的智能体强化学习系统，推理与训练采用完全异步的策略分离部署，并在真实沙箱环境中打通"采样—奖励—更新"的自进化闭环。

本文在 Qwen3.5-9B 智能体上实现并端到端验证整套系统，覆盖多个能力域的工具使用任务，证明该系统能在无人工标注下自主完成"采样—奖励—更新"的自进化闭环。
\end{abstractzh}
